\documentclass[11pt,a4paper]{article}

% ── Codificación, fuentes y lenguaje ───────────────────────────────────────────
\usepackage[utf8]{inputenc}
\usepackage[T1]{fontenc}
\usepackage[default,scale=0.95]{sourcesanspro}
\renewcommand{\familydefault}{\sfdefault}
\usepackage[spanish,es-noshorthands,es-tabla]{babel}

% ── Geometría y espaciado ─────────────────────────────────────────────────────
\usepackage[top=2.2cm, bottom=2.5cm, left=2.2cm, right=2.2cm, headheight=15pt]{geometry}
\usepackage{setspace}
\setstretch{1.18}
\usepackage{parskip}

% ── Matemáticas y unidades ────────────────────────────────────────────────────
\usepackage{amsmath}
\usepackage{amssymb}
\usepackage{siunitx}

% ── Circuitos ─────────────────────────────────────────────────────────────────
\usepackage[european, straightvoltages]{circuitikz}

% ── Tablas ────────────────────────────────────────────────────────────────────
\usepackage{booktabs}
\usepackage{tabularx}
\usepackage{array}
\usepackage{longtable}
\usepackage{colortbl}
\usepackage{enumitem}

% ── Colores, cajas e iconos ───────────────────────────────────────────────────
\usepackage[dvipsnames]{xcolor}
\usepackage{tcolorbox}
\tcbuselibrary{skins, breakable}
\usepackage{fontawesome5}
\usepackage{graphicx}

% ── Listados de código (dark-mode infográfico) ────────────────────────────────
\usepackage{listings}

% ── Secciones con barra de color (infografía) ────────────────────────────────
\usepackage{titlesec}
\titleformat{\section}
  {\normalfont\LARGE\bfseries\color{azulNoche}}
  {\colorbox{cyanNeon}{\color{fondoOscuro}\thesection}}{0.7em}{}
  [\vspace{2pt}\color{cyanNeon}\rule{\linewidth}{1.8pt}]
\titlespacing*{\section}{0pt}{24pt}{10pt}
\titleformat{\subsection}
  {\normalfont\large\bfseries\color{azulNoche}}
  {\textcolor{cyanNeon}{\thesubsection}}{0.7em}{}
  [\vspace{1pt}\color{grisLinea}\rule{\linewidth}{0.6pt}]
\titlespacing*{\subsection}{0pt}{14pt}{6pt}
\titleformat{\subsubsection}
  {\normalfont\normalsize\bfseries\color{azulNoche}}
  {\textcolor{cyanNeon}{\thesubsubsection}}{0.7em}{}
\titlespacing*{\subsubsection}{0pt}{10pt}{4pt}

% ── Cabeceras y pies de página ────────────────────────────────────────────────
\usepackage{fancyhdr}
\usepackage{lastpage}
\pagestyle{fancy}
\fancyhf{}
\renewcommand{\headrulewidth}{0pt}
\renewcommand{\footrulewidth}{0pt}
\fancyfoot[C]{%
  \begin{minipage}{\linewidth}
    \vspace{2pt}
    \color{cyanNeon}\rule{\linewidth}{1.2pt}\\[2pt]
    \centering\color{grisTexto}\footnotesize
    Página \thepage\ de \pageref{LastPage}
  \end{minipage}%
}

% ── Hiperenlaces ──────────────────────────────────────────────────────────────
\usepackage[colorlinks=true, linkcolor=azulNoche, urlcolor=cyanNeon]{hyperref}
\hypersetup{
    pdftitle={Veredicto de Hardware: HP EliteDesk 805 G6 Mini},
}
\sloppy
\emergencystretch 3em

% ── Paleta de colores — tecnológico dark-mode ──────────────────────────────────
\definecolor{fondoOscuro}{HTML}{0D1B2A}     % Dark navy — fondo banda título
\definecolor{verdeTurquesa}{HTML}{004D40}    % Verde turquesa — bandas de marca
\definecolor{azulNoche}{HTML}{1B3A6B}       % Midnight blue — secciones, marcos
\definecolor{azulMedio}{HTML}{2A5298}       % Azul medio — variante de acento
\definecolor{cyanNeon}{HTML}{00C8E0}        % Cyan eléctrico — acentos primarios
\definecolor{verdeSignal}{HTML}{00B887}     % Verde señal — fortalezas, éxito
\definecolor{naranjaVivo}{HTML}{FF7043}     % Naranja vivo — mejoras, advertencias
\definecolor{rojoAlerta}{HTML}{E53935}      % Rojo alerta — errores
\definecolor{grisPapel}{HTML}{F4F6FA}       % Fondo tarjetas claras
\definecolor{grisLinea}{HTML}{C8D0E0}       % Bordes sutiles
\definecolor{grisTexto}{HTML}{6B7A99}       % Texto secundario
\definecolor{amarilloNota}{HTML}{FFB300}    % Notas / advertencias doradas

% Colores específicos para bloques de código (dark-mode)
\definecolor{codigoFondo}{HTML}{1E2D3D}      % Fondo del bloque de código
\definecolor{codigoComentario}{HTML}{7A8FA6} % Comentarios
\definecolor{codigoCadena}{HTML}{FF9D5C}     % Cadenas / strings
\definecolor{codigoKeyword}{HTML}{00C8E0}    % Palabras clave
\definecolor{codigoNumero}{HTML}{00E5A0}     % Números

% ── Banda de título: portada infográfica ──────────────────────────────────────
%   Diseño en 2 capas: banda de color (por defecto fondo oscuro) + franja de
%   acento cyan abajo. Uso: \bandaTitulo[color]{título}{subtítulo}.
%   El icono se puede cambiar con \renewcommand{\iconoBanda}{...} (FontAwesome).
\newcommand{\iconoBanda}{microchip}
\newcommand{\bandaTitulo}[3][fondoOscuro]{%
  \begin{tcolorbox}[
    enhanced,
    colback=#1, colframe=#1,
    boxrule=0pt, arc=8pt, outer arc=8pt,
    width=\linewidth,
    top=18pt, bottom=6pt, left=14pt, right=14pt,
    borderline south={3.5pt}{0pt}{cyanNeon},
  ]
    \begin{minipage}[c]{0.07\linewidth}
      \centering
      \textcolor{cyanNeon}{\fontsize{30}{30}\selectfont\faIcon{\iconoBanda}}
    \end{minipage}%
    \hspace{8pt}%
    \begin{minipage}[c]{0.88\linewidth}
      {\fontsize{20}{24}\selectfont\bfseries\color{white}#2}\par\vspace{5pt}%
      \textcolor{cyanNeon!80!white}{\small\faIcon{angle-right}~#3}%
    \end{minipage}
    \vspace{4pt}
  \end{tcolorbox}%
  \vspace{8pt}%
}

% ── Tarjetas de datos (infografía) ────────────────────────────────────────────
\tcbset{
  tarjetaDato/.style={
    enhanced, breakable,
    arc=6pt, outer arc=6pt,
    colback=grisPapel, colframe=azulNoche,
    boxrule=1pt,
    fonttitle=\bfseries\small, coltitle=white,
    attach boxed title to top left={yshift=-3mm, xshift=6mm},
    boxed title style={
      colback=azulNoche, colframe=azulNoche,
      arc=4pt, boxrule=0pt,
    },
    drop fuzzy shadow=azulNoche!30!white,
    top=8pt, bottom=8pt, left=10pt, right=10pt,
  },
  tarjetaIcono/.style={
    enhanced, breakable,
    arc=6pt, outer arc=6pt,
    colback=white, colframe=grisLinea, boxrule=0.8pt,
    fonttitle=\bfseries\small, coltitle=azulNoche,
    borderline west={3pt}{0pt}{cyanNeon},
    drop fuzzy shadow=azulNoche!25!white,
    top=6pt, bottom=6pt, left=10pt, right=10pt,
  },
  cajaTitulo/.style={
    enhanced, breakable,
    arc=6pt, outer arc=6pt,
    colback=azulNoche!10!grisPapel, colframe=azulNoche,
    boxrule=1pt,
    fonttitle=\bfseries, coltitle=white,
    attach boxed title to top left={yshift=-3mm, xshift=6mm},
    boxed title style={
      colback=azulNoche, colframe=azulNoche,
      arc=4pt, boxrule=0pt,
    },
    drop fuzzy shadow=azulNoche!25!white,
    top=8pt, bottom=8pt, left=10pt, right=10pt,
  },
  cajaContenido/.style={
    enhanced, breakable,
    arc=5pt, outer arc=5pt,
    colback=grisPapel, colframe=grisLinea,
    boxrule=0.8pt,
    fonttitle=\bfseries\small, coltitle=azulNoche,
    top=6pt, bottom=6pt, left=10pt, right=10pt,
  },
  cajaFortaleza/.style={
    enhanced, breakable,
    arc=6pt, outer arc=6pt,
    colback=verdeSignal!8!white, colframe=verdeSignal,
    boxrule=1pt,
    borderline west={4pt}{0pt}{verdeSignal},
    fonttitle=\bfseries\small, coltitle=verdeSignal!70!black,
    drop fuzzy shadow=verdeSignal!20!white,
    top=6pt, bottom=6pt, left=10pt, right=10pt,
  },
  cajaMejora/.style={
    enhanced, breakable,
    arc=6pt, outer arc=6pt,
    colback=naranjaVivo!8!white, colframe=naranjaVivo,
    boxrule=1pt,
    borderline west={4pt}{0pt}{naranjaVivo},
    fonttitle=\bfseries\small, coltitle=naranjaVivo!80!black,
    drop fuzzy shadow=naranjaVivo!20!white,
    top=6pt, bottom=6pt, left=10pt, right=10pt,
  },
  cajaRecomendacion/.style={
    enhanced, breakable,
    arc=6pt, outer arc=6pt,
    colback=cyanNeon!6!white, colframe=cyanNeon!60!azulNoche,
    boxrule=1pt,
    borderline west={4pt}{0pt}{cyanNeon},
    fonttitle=\bfseries\small, coltitle=azulNoche,
    drop fuzzy shadow=azulNoche!20!white,
    top=6pt, bottom=6pt, left=10pt, right=10pt,
  },
}

% ── Ícono + texto (encabezados de sección y elementos) ────────────────────────
\newcommand{\iconotexto}[2]{\textcolor{cyanNeon}{\faIcon{#1}}~#2}

% ── Listados de código: estilo dark + auto-envoltura ─────────────────────────
\lstdefinestyle{estiloCodigo}{
    backgroundcolor=\color{codigoFondo},
    basicstyle=\ttfamily\small\color{white},
    commentstyle=\color{codigoComentario}\itshape,
    stringstyle=\color{codigoCadena},
    keywordstyle=\color{codigoKeyword}\bfseries,
    numberstyle=\tiny\color{codigoComentario},
    numbers=left,
    numbersep=10pt,
    showstringspaces=false,
    breaklines=true,
    breakatwhitespace=false,
    tabsize=4,
    frame=none,
    captionpos=b,
    aboveskip=6pt,
    belowskip=4pt,
    xleftmargin=14pt,
    framexleftmargin=14pt,
    literate=
      {á}{{\'a}}1 {é}{{\'e}}1 {í}{{\'i}}1 {ó}{{\'o}}1 {ú}{{\'u}}1
      {Á}{{\'A}}1 {É}{{\'E}}1 {Í}{{\'I}}1 {Ó}{{\'O}}1 {Ú}{{\'U}}1
      {ñ}{{\~n}}1 {Ñ}{{\~N}}1 {ü}{{\"u}}1 {Ü}{{\"U}}1
      {¿}{{?`}}1 {¡}{{!`}}1
}
\lstdefinestyle{estiloPython}{
    style=estiloCodigo,
    language=Python,
}
\lstset{style=estiloCodigo}
\BeforeBeginEnvironment{lstlisting}{%
  \begin{tcolorbox}[
    enhanced, breakable,
    arc=6pt, outer arc=6pt,
    colback=codigoFondo, colframe=azulNoche,
    boxrule=1pt,
    drop fuzzy shadow=azulNoche!25!white,
    top=2pt, bottom=2pt, left=0pt, right=0pt,
  ]%
}
\AfterEndEnvironment{lstlisting}{\end{tcolorbox}}

% ── Cajas de contenido temáticas ──────────────────────────────────────────────
\newtcolorbox{cajaRecuerda}{
    enhanced, breakable,
    arc=6pt, outer arc=6pt,
    colback=azulNoche!10!grisPapel,
    colframe=azulNoche, boxrule=1pt,
    borderline west={4pt}{0pt}{cyanNeon},
    fonttitle=\bfseries\small\color{white},
    title={\faIcon{info-circle}~Recuerda},
    attach boxed title to top left={yshift=-3mm, xshift=6mm},
    boxed title style={colback=azulNoche, arc=4pt, boxrule=0pt},
    drop fuzzy shadow=azulNoche!25!white,
    left=10pt, right=10pt, top=8pt, bottom=8pt,
}

\newtcolorbox{cajaConcepto}{
    enhanced, breakable,
    arc=6pt, outer arc=6pt,
    colback=verdeSignal!8!white,
    colframe=verdeSignal, boxrule=1pt,
    borderline west={4pt}{0pt}{verdeSignal},
    fonttitle=\bfseries\small\color{white},
    title={\faIcon{lightbulb}~Concepto clave},
    attach boxed title to top left={yshift=-3mm, xshift=6mm},
    boxed title style={colback=verdeSignal!80!black, arc=4pt, boxrule=0pt},
    drop fuzzy shadow=verdeSignal!20!white,
    left=10pt, right=10pt, top=8pt, bottom=8pt,
}

\newtcolorbox{cajaEjemplo}{
    enhanced, breakable,
    arc=6pt, outer arc=6pt,
    colback=cyanNeon!5!white,
    colframe=cyanNeon!70!azulNoche, boxrule=1pt,
    borderline west={4pt}{0pt}{cyanNeon},
    fonttitle=\bfseries\small\color{fondoOscuro},
    title={\faIcon{code}~Ejemplo},
    attach boxed title to top left={yshift=-3mm, xshift=6mm},
    boxed title style={colback=cyanNeon, arc=4pt, boxrule=0pt},
    drop fuzzy shadow=azulNoche!20!white,
    left=10pt, right=10pt, top=8pt, bottom=8pt,
}

\newtcolorbox{cajaEjercicio}{
    enhanced, breakable,
    arc=6pt, outer arc=6pt,
    colback=azulMedio!7!white,
    colframe=azulMedio, boxrule=1pt,
    borderline west={4pt}{0pt}{azulMedio},
    fonttitle=\bfseries\small\color{white},
    title={\faIcon{pencil-alt}~Ejercicio},
    attach boxed title to top left={yshift=-3mm, xshift=6mm},
    boxed title style={colback=azulMedio, arc=4pt, boxrule=0pt},
    drop fuzzy shadow=azulNoche!20!white,
    left=10pt, right=10pt, top=8pt, bottom=8pt,
}

\newtcolorbox{cajaReto}{
    enhanced, breakable,
    arc=6pt, outer arc=6pt,
    colback=naranjaVivo!6!white,
    colframe=naranjaVivo, boxrule=1pt,
    borderline west={4pt}{0pt}{naranjaVivo},
    fonttitle=\bfseries\small\color{white},
    title={\faIcon{fire}~Reto},
    attach boxed title to top left={yshift=-3mm, xshift=6mm},
    boxed title style={colback=naranjaVivo, arc=4pt, boxrule=0pt},
    drop fuzzy shadow=naranjaVivo!20!white,
    left=10pt, right=10pt, top=8pt, bottom=8pt,
}

% ── Indicadores de dificultad ─────────────────────────────────────────────────
\newcommand{\facil}{\textcolor{verdeSignal}{\faIcon{circle}\,\textbf{Fácil}}}
\newcommand{\intermedio}{\textcolor{naranjaVivo}{\faIcon{circle}\,\textbf{Intermedio}}}
\newcommand{\dificil}{\textcolor{rojoAlerta}{\faIcon{circle}\,\textbf{Difícil}}}

% ─────────────────────────────────────────────────────────────────────────────
% Cabecera específica de este documento
\usepackage{url}
\usepackage{needspace}
\hypersetup{
    pdftitle={Graphify: evaluación de incorporación en ELECTRONICA},
}
\fancyhead[L]{\small\color{azulNoche}\textbf{Agente IA --- Evaluación de incorporación}}
\fancyhead[R]{\small\color{grisTexto}\itshape Graphify v0.9.85}
\renewcommand{\headrulewidth}{0.6pt}
\renewcommand{\headrule}{\color{cyanNeon}\hrule width\headwidth height 0.6pt}

\renewcommand{\iconoBanda}{project-diagram}

\begin{document}

% ══ PORTADA INFográfica ══════════════════════════════════════════════════════
\bandaTitulo{Graphify: evaluación de incorporación}{Grafo de conocimiento determinista para ELECTRONICA --- guía para el HP EliteDesk 805 G6}

\begin{center}
  \vspace{2pt}
  {\large \textbf{Autor:} Prof. CERO \hfill \textbf{Fecha:} \today} \\[6pt]
  \color{cyanNeon}\rule{\linewidth}{0.4pt}
\end{center}

\tableofcontents
\newpage

% ══ 1. RESUMEN EJECUTIVO ═════════════════════════════════════════════════════
\section[Resumen Ejecutivo]{\iconotexto{info-circle}{Resumen Ejecutivo}}

Se evaluó \textbf{Graphify} (\texttt{Graphify-Labs/graphify}, versión \textbf{v0.9.85},
publicada el 10 de octubre de 2026) como herramienta de representación y consulta del
código del workspace \texttt{ELECTRONICA}. Graphify convierte un directorio de código,
documentación, esquemas SQL, configuraciones y PDFs en un \textbf{grafo de conocimiento
consultable} (\texttt{graph.json}) construido de forma \textbf{determinista} a partir de
árboles sintácticos (\emph{tree-sitter}), \textbf{sin LLM y sin base de datos vectorial}.

El veredicto es \textbf{adoptable como herramienta de solo-código, tras un piloto
medido}. Su principio de diseño ---un script determinista que produce un artefacto
reproducible--- es exactamente el que sostiene la arquitectura de 3 capas del proyecto.
La migración a la \textbf{HP EliteDesk 805 G6} (16\,GB de RAM) elimina además el único
obstáculo de recursos que hoy lo bloquea en el equipo Celeron (3{,}63\,GB).

\begin{tcolorbox}[cajaFortaleza, title={\faIcon{thumbs-up}~Veredicto rápido: SÍ, adoptable tras un piloto}]
Graphify aporta una capa que el proyecto \textbf{hoy no tiene}: las relaciones
\texttt{calls}/\texttt{imports}/\texttt{inherits} resueltas \emph{entre ficheros} como un
grafo tipado y verificable. \textbf{No sustituye} al RAG actual
(\texttt{rag\_system.py}): son complementarios. La única regla no negociable es que su
\emph{pase semántico} (docs/PDF/imágenes, que sí usa un LLM) no debe romper el
determinismo del proyecto.
\end{tcolorbox}

\begin{cajaRecuerda}
\textbf{Guardarraíl determinista.} Adoptar Graphify como \emph{herramienta externa} no
viola las 3 capas: no introduce lógica de decisión en el chat ni un flujo nuevo. Sí lo
violaría (a) resolver su pase semántico en prosa de chat, o (b) crear un
\texttt{flujo\_*} que lo use sin su directiva y su script. Cualquier integración real
exige las tres capas: \texttt{directives/*.yaml} + orquestador + \texttt{execution/*.py}.
\end{cajaRecuerda}

% ══ 2. QUÉ ES GRAPHIFY ══
\section[Qué es Graphify]{\iconotexto{sitemap}{Qué es Graphify}}

Graphify es una \emph{skill} para asistentes de IA respaldada por una biblioteca Python.
La \emph{skill} orquesta la biblioteca; la biblioteca puede usarse sola. El paquete de
PyPI se llama \texttt{graphifyy} (doble ``y''), mientras que el comando de consola es
\texttt{graphify}.

\begin{tcolorbox}[tarjetaDato, title={\faIcon{project-diagram}~Ficha de la herramienta}]
\begin{tabular}{ll}
\toprule
\textbf{Repositorio} & \texttt{Graphify-Labs/graphify} \\
\textbf{Versión evaluada} & v0.9.85 (10 de octubre de 2026) \\
\textbf{Licencia} & Apache-2.0 + MIT \\
\textbf{Lenguaje} & Python (\(\ge\) 3.10) \\
\textbf{Paquete PyPI} & \texttt{graphifyy} (comando: \texttt{graphify}) \\
\textbf{Dependencias base} & \texttt{networkx}, \texttt{numpy}, \texttt{rapidfuzz} y \(\sim\)26 gramáticas \texttt{tree-sitter} \\
\textbf{Salida} & \texttt{graph.json}, \texttt{graph.html}, \texttt{GRAPH\_REPORT.md} \\
\textbf{Vectores} & \textbf{Ninguno} (no hay \emph{embeddings}) \\
\textbf{Telemetría} & Ninguna (registra solo consultas en \texttt{\textasciitilde/.cache/}) \\
\bottomrule
\end{tabular}
\end{tcolorbox}

\subsection{El pipeline determinista}

Toda la construcción del grafo para \textbf{código} es local y sin LLM:

\begin{center}
\texttt{detect()} \(\rightarrow\) \texttt{extract()} \(\rightarrow\) \texttt{build()}
\(\rightarrow\) \texttt{cluster()} \(\rightarrow\) \texttt{report} \(\rightarrow\)
\texttt{export}
\end{center}

Cada etapa vive en su propio módulo y se comunican por diccionarios y grafos
\texttt{networkx}. El código se analiza con \textbf{árboles sintácticos} (tree-sitter,
unos 40 lenguajes), \textbf{sin enviar nada fuera de la máquina}. Solo el pase semántico
sobre documentación, PDFs, imágenes y vídeo requiere un \emph{backend} LLM, y únicamente
si se le configura uno.

\begin{cajaConcepto}
\textbf{Cada arista está etiquetada.} Una conexión se marca \texttt{EXTRACTED} si es
explícita en el código fuente (un \texttt{import}, una llamada directa) e
\texttt{INFERRED} si la dedujo el resolver. Las ambiguas se marcan \texttt{AMBIGUOUS} y
se aíslan para revisión humana. Esto es lo que separa a Graphify de una caja negra: se
puede auditar qué se leyó y qué se dedujo.
\end{cajaConcepto}

\subsection{Artefactos que produce}

\needspace{6\baselineskip}
\begin{lstlisting}
graphify-out/
|-- graph.html        (grafo interactivo: abrir en el navegador)
|-- GRAPH_REPORT.md   (god nodes, comunidades, preguntas sugeridas)
`-- graph.json        (el grafo completo, consultable sin releer ficheros)
\end{lstlisting}

El grafo incluye los ``god nodes'' (los conceptos más conectados), las comunidades
detectadas por el algoritmo de Leiden y rutas entre símbolos (\texttt{path}) o
explicaciones de un concepto (\texttt{explain}).

\begin{tcolorbox}[cajaMejora, title={\faIcon{exclamation-triangle}~Aviso de deriva (\emph{churn})}]
\textbf{Cadencia de publicación muy alta.} Entre el 6 y el 10 de octubre de 2026 se
publicaron las versiones 0.9.78 a 0.9.85 (una por día, aproximadamente). El ritmo es una
señal de salud del proyecto, pero también de \textbf{deriva de esquema y de API}: por eso
el grafo lleva \texttt{schema\_version} y \texttt{graphify\_version}, y conviene fijar una
versión concreta en lugar de seguir \texttt{latest}.
\end{tcolorbox}

% ══ 3. ENCAJE CON LA ARQUITECTURA ══
\section[Encaje con ELECTRONICA]{\iconotexto{balance-scale}{Encaje con la arquitectura de 3 capas}}

\begin{tcolorbox}[tarjetaIcono, title={\faIcon{check-circle}~Principio rector compartido}]
\textbf{Electrónica} separa orquestación probabilística de ejecución determinista.
\textbf{Graphify} hace lo mismo: un script determinista (\texttt{extract}, tree-sitter)
produce un artefacto estable; el LLM solo lo consume. No son dos filosofías distintas,
es la misma aplicada al código como grafo.
\end{tcolorbox}

\begin{itemize}[leftmargin=1.4em]
    \item \textbf{Determinismo.} Los últimos ocho lanzamientos son, casi en su totalidad,
    correcciones de determinismo y portabilidad: los identificadores de nodo ya no
    filtran la ruta de \emph{checkout} ni el usuario del sistema operativo, de modo que
    \textbf{un clon produce el mismo grafo}. Es la misma garantía ``mismo input
    \(\rightarrow\) mismo output'' del enrutador.
    \item \textbf{Integración con OpenCode.} Graphify soporta explícitamente
    \texttt{--platform opencode}; el \emph{plugin} admite el contrato v1 (\texttt{server()})
    y v2 (\texttt{setup()}). Encaja con el motor rotativo del asistente y, en modo
    solo-código, \textbf{no consume} la clave de OpenRouter.
    \item \textbf{Servidor MCP opcional.} El extra \texttt{graphifyy[mcp]} expone un
    servidor (\texttt{graphify-mcp}) que se suma a los servidores MCP del proyecto. Es un
    proceso aparte, por lo que no colisiona con el pin de \texttt{mcp\,<\,2} del resto.
    \item \textbf{Seguridad.} Todas las entradas externas pasan por
    \texttt{graphify/security.py}: validación de URL (solo http/https, con bloqueo de
    redirecciones \texttt{file://}), límite de tamaño y tiempo, y rutas de grafo que deben
    resolverse \emph{dentro} de \texttt{graphify-out/}.
\end{itemize}

% ══ 4. BENEFICIOS ══
\section[Beneficios]{\iconotexto{chart-line}{Beneficios concretos}}

\begin{itemize}[leftmargin=1.4em]
    \item \textbf{Consultar en vez de \emph{grep}.} \texttt{graphify query},
    \texttt{graphify path} y \texttt{graphify explain} sustituyen la lectura fichero a
    fichero, lo que reduce los tokens que consume el orquestador.
    \item \textbf{Relaciones entre ficheros.} Resuelve \texttt{calls}, \texttt{imports},
    \texttt{inherits} y \texttt{mixes\_in} entre módulos, algo que el RAG actual
    \textbf{no} hace.
    \item \textbf{Verificable.} Las etiquetas \texttt{EXTRACTED}/\texttt{INFERRED} y la
    resolución \emph{fail-closed} ante ambigüedad hacen el resultado auditable.
    \item \textbf{Racional y referencias.} Los comentarios \texttt{\# NOTE:} / \texttt{\# WHY:}
    y las citas a ADR/RFC se vuelven nodos de primera clase enlazados al código.
    \item \textbf{Coste cero en solo-código.} \texttt{extract --code-only} corre
    completamente offline y no gasta créditos.
\end{itemize}

% ══ 5. RIESGOS Y COSTES ══
\section[Riesgos y Costes]{\iconotexto{shield-alt}{Riesgos y costes}}

\begin{longtable}{p{3.3cm} p{8.4cm} p{2.0cm}}
\toprule
\textbf{Riesgo} & \textbf{Detalle} & \textbf{Severidad} \\
\midrule
\endhead
\bottomrule
\endlastfoot
Recursos & En el equipo Celeron había \textbf{641\,MB} de RAM disponibles y el repo pesa
\textbf{1{,}4\,GB} (2105 ficheros). Con la HP EliteDesk (16\,GB) el riesgo \textbf{baja a
residual} para solo-código; el paso semántico de PDFs/imágenes sigue siendo el tramo
pesado. & Media \(\rightarrow\) Baja \\
\addlinespace
Dependencias & \texttt{uv tool install} crea un entorno aislado (no toca
\texttt{elect\_env}), pero descarga \(\sim\)26 gramáticas de tree-sitter; el extra
\texttt{leiden} arrastra \texttt{graspologic} (scipy/sklearn/numba) en Python\,<\,3.13.
& Media \\
\addlinespace
Edita el árbol de trabajo & Escribe \texttt{graphify-out/} en el directorio actual y,
con \texttt{graphify hook install}, \textbf{reconstruye en cada commit}. & Media \\
\addlinespace
Edita \texttt{AGENTS.md} & En OpenCode Graphify es plataforma de fichero de
instrucciones: \texttt{install} inserta guía en \texttt{AGENTS.md}, que en este repo es
\textbf{la fuente de verdad}. Un merge automático puede dañarla. & \textbf{Alta} \\
\addlinespace
Créditos / guardarraíl & El pase semántico manda docs/PDFs a un \emph{backend} LLM. Usarlo
\emph{inline} violaría el guardarraíl. Debe ir por \texttt{execution/enrutador.py} o
quedarse en \texttt{--code-only}. & Media \\
\addlinespace
Deriva (\emph{churn}) & Cadencia casi diaria y marcado enfoque a la nube. Conviene fijar
una versión. & Media \\
\addlinespace
Disco de 256\,GB & El SSD del equipo elegido es modesto (ya al 91\,\% en el equipo actual);
mantener \texttt{chroma\_db}, \texttt{datasets} y modelos fuera del disco de sistema. &
Media \\
\end{longtable}

% ══ 6. COMPARACIÓN CON EL RAG ACTUAL ══
\section[Comparación con el RAG]{\iconotexto{search}{Comparación con el RAG actual}}

\begin{longtable}{p{3.4cm} p{5.1cm} p{5.1cm}}
\toprule
\textbf{Aspecto} & \textbf{Graphify} & \textbf{rag\_system.py (ChromaDB)} \\
\midrule
\endhead
\bottomrule
\endlastfoot
Unidad & Grafo tipado: nodos = conceptos, aristas explicadas & \emph{Embeddings} y similitud \\
Determinismo & Sí (para código) & No (vectores aproximados) \\
Almacén vectorial & \textbf{No} & Sí (\texttt{chroma\_db/}) \\
Cubre código & Sí (unos 40 lenguajes) & No (docs, \texttt{.tex}, \texttt{.pdf}) \\
Coste por construcción & 0 créditos en \texttt{--code-only} & LLM por consulta \\
Relación & Estructura (``qué llama a qué'') & Contenido (``qué dice qué'') \\
\end{longtable}

\begin{cajaConcepto}
\textbf{No es un sustituto, es una capa complementaria.} Graphify aporta la
\emph{estructura} del código; el RAG aporta el \emph{contenido} de la documentación.
Conviven sin solaparse: una pregunta ``¿qué llama a \texttt{calcular\_nota}?'' es de
Graphify; ``¿qué dice la rúbrica sobre el montaje?'' es del RAG.
\end{cajaConcepto}

% ══ 7. VEREDICTO ══
\section[Veredicto]{\iconotexto{gavel}{Veredicto}}

\begin{tcolorbox}[cajaRecomendacion, title={\faIcon{clipboard-check}~Recomendación}]
\textbf{Adoptar Graphify como herramienta de solo-código, en tres fases y sin tocar
todavía la arquitectura del repo.} El cuello de botella real del equipo Celeron era la
RAM y el volumen, no el software ---el mismo patrón que ya se documentó para el MCP de
FreeCAD---. La migración a la HP EliteDesk resuelve la RAM; el volumen del repo y el
disco de 256\,GB siguen siendo el punto a medir.
\end{tcolorbox}

\begin{itemize}[leftmargin=1.4em]
    \item \textbf{Fase 0 --- Alcance.} Decidir si el piloto cubre solo los lenguajes vivos
    (\texttt{execution/}, \texttt{.opencode/}, los \texttt{flujo\_*.py}) o el repo completo.
    \item \textbf{Fase 1 --- Piloto aislado.} Instalar, construir el grafo y medir, sin
    modificar ningún fichero del repo.
    \item \textbf{Fase 2 --- Instalación por proyecto} (\texttt{--project}), revisando y
    \textbf{mergeando a mano} el diff de \texttt{AGENTS.md}.
    \item \textbf{Fase 3 --- Integración real} (si se quiere ir más allá): nueva pareja
    \texttt{directives/graphify\_code.yaml} + \texttt{execution/graphify\_build.py}.
\end{itemize}

% ══ 8. GUÍA DE IMPLEMENTACIÓN EN LA HP ELITEDESK 805 G6 ══
\section[Implementación en la HP]{\iconotexto{laptop}{Guía de implementación en la HP EliteDesk 805 G6}}

Esta sección es la operativa: qué ejecutar, en qué orden, y cómo comprobar que salió bien.

\subsection{Contexto de hardware}

\begin{tcolorbox}[tarjetaDato, title={\faIcon{microchip}~HP EliteDesk 805 G6 Mini (equipo destino)}]
\begin{tabular}{ll}
\toprule
\textbf{Procesador} & AMD Ryzen 5 PRO 4650GE (Zen 2, 6C/12T, \SI{35}{\watt}) \\
\textbf{Instrucciones SIMD} & AVX2 + FMA3 (clave frente al Celeron, sin AVX) \\
\textbf{Memoria} & 16\,GB DDR4 (frente a 3{,}63\,GB del equipo actual) \\
\textbf{Almacenamiento} & SSD 256\,GB (ampliable: 1\(\times\) M.2 2280 y bahía 2{,}5'' SATA) \\
\textbf{Sistema operativo} & Windows 11 Pro \(\rightarrow\) se instalará Linux \\
\bottomrule
\end{tabular}
\end{tcolorbox}

\begin{cajaRecuerda}
El cómputo pesado sigue en \textbf{Google Colab}. Esta máquina habilita lo que el
proyecto permite en local: inferencia, prototipos, validación y orquestación. Graphify
cae del lado de la \emph{orquestación}, no del entrenamiento.
\end{cajaRecuerda}

\subsection{Prerrequisitos}

\begin{longtable}{p{4.4cm} p{3.0cm} p{6.3cm}}
\toprule
\textbf{Requisito} & \textbf{Mínimo} & \textbf{Comprobación / instalación} \\
\midrule
\endhead
\bottomrule
\endlastfoot
Python & 3.10+ & \texttt{python3 --version} (el \texttt{elect\_env} es 3.10) \\
\texttt{uv} (recomendado) & cualquiera & \path{curl -LsSf https://astral.sh/uv/install.sh | sh} \\
\texttt{pipx} (alternativa) & cualquiera & \texttt{pip install pipx} \\
\bottomrule
\end{longtable}

En esta máquina de trabajo ya hay \texttt{uvx} 0.11.11 y Python 3.10.20, así que la ruta
de \texttt{uv} es la natural.

\subsection{Instalación (entorno aislado)}

\begin{lstlisting}
# 1. Instalar el CLI en un entorno aislado (no toca el entorno elect_env)
uv tool install graphifyy

# 2. Registrar la skill en el asistente OpenCode
graphify install --platform opencode
\end{lstlisting}

\begin{tcolorbox}[cajaMejora, title={\faIcon{exclamation-triangle}~No dejar que reescriba \texttt{AGENTS.md}}]
\textbf{No} ejecutar \texttt{graphify install} sin plataforma: por diseño apunta a
Claude Code. La forma correcta aquí es \texttt{graphify install --platform opencode}. Y
antes de aceptar cualquier escritura sobre \texttt{AGENTS.md}, \textbf{revisar el diff}
y fusionarlo a mano; nunca dejar que lo sobreescriba.
\end{tcolorbox}

Extras opcionales (instalar solo lo que se use):

\begin{lstlisting}
uv tool install "graphifyy[leiden]"   # deteccion de comunidades (Leiden)
uv tool install "graphifyy[mcp]"      # servidor MCP (graphify-mcp)
uv tool install "graphifyy[pdf]"      # extraccion de PDF (pase semantico)
uv tool install "graphifyy[all]"      # todo lo anterior y mas
\end{lstlisting}

\subsection{Primera construcción del grafo}

Para arrancar sin gastar créditos, indexar \textbf{solo el código} (offline):

\begin{lstlisting}
cd ~/ELECTRONICA
graphify extract --code-only .      # construye graphify-out/ sin LLM
\end{lstlisting}

Dentro del asistente, la skill ofrece el atajo:

\begin{lstlisting}
/graphify .                         # construir el grafo de la carpeta actual
/graphify ./docs --update           # re-extraer solo los ficheros cambiados
/graphify . --no-viz                # omitir el HTML, solo informe + JSON
\end{lstlisting}

Consultas habituales una vez construido el grafo:

\begin{lstlisting}
graphify query "que conecta el enrutador con el cliente LLM?"
graphify path "enrutador.py" "llm_client.py"
graphify explain "calcular_nota"
\end{lstlisting}

\subsection{Verificación de determinismo}

El valor del proyecto está en la reproducibilidad: \textbf{dos construcciones deben
producir el mismo \texttt{graph.json}}.

\begin{lstlisting}
graphify extract --code-only . && sha256sum graphify-out/graph.json > /tmp/g1.sha
graphify extract --code-only . && sha256sum graphify-out/graph.json > /tmp/g2.sha
diff /tmp/g1.sha /tmp/g2.sha && echo "DETERMINISTA: hashes identicos"
\end{lstlisting}

Si los hashes difieren, \textbf{detenerse y diagnosticar} antes de dar el piloto por
bueno: un grafo no reproducible contradice el guardarraíl del proyecto.

\subsection{Integración opcional}

\begin{itemize}[leftmargin=1.4em]
    \item \textbf{Servidor MCP.} Con el extra \texttt{[mcp]}, \texttt{graphify-mcp}
    expone el grafo como herramientas MCP, sumándose a los servidores ya existentes.
    \item \textbf{Hook de git.} \texttt{graphify hook install} reconstruye el grafo en cada
    commit y \emph{checkout}. Útil, pero añade coste por commit: evaluar antes de dejarlo
    permanente.
    \item \textbf{Pase semántico.} Si en algún momento se indexan docs/PDFs, la llamada al
    LLM debe pasar por \texttt{execution/enrutador.py} (\texttt{--task destilacion}), nunca
    decidirse en el chat.
    \item \textbf{Control de versiones.} Decidir conscientemente si \texttt{graphify-out/}
    se ignora o se comparte; no dejarlo al azar.
\end{itemize}

\subsection{Convivencia con la arquitectura de 3 capas}

\begin{cajaReto}
Si la adopción pasa de \emph{herramienta} a \emph{flujo}, hace falta su tríada
completa: \texttt{directives/graphify\_code.yaml} (el \emph{qué}),
un orquestador que valide entradas/salidas, y \texttt{execution/graphify\_build.py}
como capa 3 determinista (envoltorio que llama al CLI, valida el JSON y respeta el
\emph{retry budget} de 3 intentos). Nunca un orquestador sin directiva ni script.
\end{cajaReto}

% ══ 9. CRITERIOS DE ACEPTACIÓN ══
\section[Criterios de aceptación]{\iconotexto{clipboard-check}{Criterios de aceptación}}

El piloto se considera superado cuando se cumplan \textbf{todos} estos puntos:

\begin{itemize}[leftmargin=1.4em]
    \item \texttt{uv tool install graphifyy} deja el comando \texttt{graphify}
    disponible sin romper \texttt{elect\_env}.
    \item \texttt{graphify extract --code-only .} termina con código 0 y no consume
    créditos de OpenRouter.
    \item Dos construcciones consecutivas producen el \textbf{mismo \texttt{graph.json}}
    (hashes idénticos).
    \item \texttt{GRAPH\_REPORT.md} identifica \emph{god nodes} y comunidades coherentes
    con la estructura real del repo.
    \item Ningún fichero del repo queda modificado salvo, si se decide, el \emph{merge}
    manual y revisado de \texttt{AGENTS.md}.
    \item Se anotan tiempos y consumo de RAM por construcción, para comparar el antes
    (Celeron) y el después (HP EliteDesk).
\end{itemize}

\begin{cajaRecuerda}
Al recibir la HP EliteDesk, antes de nada: confirmar la CPU (\textbf{4650GE}, no 4350GE),
ejecutar \texttt{execution/env\_diagnostic.py} (confirmar AVX2) y
\texttt{flujo\_auditar\_sistema.py} (veredicto de RAM y servicios). Esos dos pasos ya
están anotados como criterio de aceptación del documento de hardware.
\end{cajaRecuerda}

% ══ 10. REFERENCIAS ══
\section[Referencias]{\iconotexto{book}{Referencias}}

\begin{itemize}[leftmargin=1.4em]
    \item \textbf{Repositorio:} \url{https://github.com/Graphify-Labs/graphify} ---
    \texttt{README.md}, \texttt{ARCHITECTURE.md}, \texttt{SECURITY.md},
    \texttt{pyproject.toml} y notas de versión v0.9.76--v0.9.85.
    \item \textbf{Documentación:} \url{https://docs.graphify.com}
    (guía de inicio, referencia de CLI, configuración, entradas soportadas).
    \item \textbf{Documento de hardware hermano:}
    \texttt{docs/PC\_PARA\_IA/HP\_ELITEDESK\_VEREDICTO.tex} --- veredicto del equipo
    destino.
    \item \textbf{Sesión de continuidad:}
    \texttt{Sessions/2026-10-10\_veredicto\_hardware\_hp\_elitedesk.md} y
    \texttt{Sessions/2026-10-11\_graphify\_evaluacion\_incorporacion.md}.
\end{itemize}

\end{document}

